%%%%%%%%%%%%%%%%%%%%%%%%%%%%%%%%%%%%%%%%%%%%%%%%%%%%%%%%%%%%%%%%%%%%%%%%%%%%%%%%
%%  EINLAGE 1: ARBEITSPLAN          Zwischenabgabe: 28.10.2026 (obligatorisch)
%%  Struktur: Beschreibung (1/4–1/2 S.) + Reflexion (1–2 S.)
%%  Die Einlage selbst (der Plan) steht im Anhang: 1_arbeitsplan_anhang.tex
%%%%%%%%%%%%%%%%%%%%%%%%%%%%%%%%%%%%%%%%%%%%%%%%%%%%%%%%%%%%%%%%%%%%%%%%%%%%%%%%
\section{Einlage: Arbeitsplan}
\label{sec:arbeitsplan}

\begin{hinweise}
Diese Einlage ist Ihr Arbeitsplan. Der Arbeitsplan ist die \emph{einzige} Einlage, die mit der Partnerin/dem
Partner identisch sein darf (\enquote{Ihr gemeinsamer Arbeitsplan}). Hinweise zur Erstellung (Marti): nicht nur
eine grobe Liste, sondern feine Teilaufgaben mit zeitlichem Ablauf, sichtbar machen, was parallel läuft;
Vorbereitung (\enquote{Einkaufen}) und Nachbereitung (\enquote{Abwasch}) nicht vergessen -- z.\,B. als
Excel-Tabelle oder in ProjectLibre, oder mit dem Gantt-Gerüst im Anhang.

\smallskip\emph{Zusätzlich aus der Moodle-Aufgabe (Frist 28.\,10.\,2026, 23:59\,Uhr):}
\begin{itemize}\itemsep0pt
\item Den Plan mit den Betreuern besprechen -- deren Rückmeldung in der Reflexion aufgreifen und als
      Gespräch zitieren (Vorlage \texttt{VORLAGE-Gespraech} in \texttt{literatur.bib}).
\item Gleicher Arbeitsplan für beide, aber die Reflexion schreibt jeder \emph{einzeln und unabhängig} von
      \Partner.
\item Reflexion ist keine Nacherzählung; Beschreibung $1/4$--$1/2$ Seite, Reflexion 1--2 Seiten
      (ohne Titel, Abbildungen, Tabellen und Legenden).
\item Arbeitsplan und Reflexion als \emph{ein PDF} hochladen (\texttt{Manevski\_Arbeitsplan.pdf}) --
      danach den \textbf{Abgabeknopf} drücken!
\end{itemize}
\end{hinweise}

\subsection{Beschreibung}
\begin{hinweise}
Die Beschreibung gibt alle notwendigen Informationen an, die die Einlage beschreiben, also zum Beispiel den
Fundort, die Literaturstelle und eine kurze Inhaltsübersicht, aber keine Nacherzählung ($1/4$ Seite, max.
$1/2$ Seite!). Leitfaden: Um welche Art von Einlage handelt es sich? Woher stammt sie? Welches sind die
wesentlichen Aussagen?
\end{hinweise}
\begin{hinweise}
\emph{Deine Gliederung aus dem Fragebogen (ca.\ $1/3$ Seite, 5--6 Sätze):}
\begin{enumerate}\itemsep0pt
\item \emph{Was:} gemeinsamer Arbeitsplan mit \Partner, in Excel als Gantt-Diagramm, Wochenraster
      KW\,42/2026 bis KW\,11/2027, im Anhang~\ref{anh:arbeitsplan}.
\item \emph{Woher:} am 22.\,10.\,2026 fertiggestellt und am selben Tag mit Prof.\ Marti besprochen
      \cite{Marti2026Gespraech}; Baumstruktur je parallel laufendem Thema nach dem Beispiel
      \enquote{Mittagessen kochen} \cite{MoodleArbeitsplan2026}; alles gemeinsam erstellt.
\item \emph{Aufbau:} sieben Blöcke (Theorie \& Literatur, Materialien, Versuchsaufbau, Messungen,
      Auswertung, Vortrag, Portfolio); jede Aufgabe mit Person, Start, Ende, Status; Prozentangabe je
      Block; alle Abgaben eingetragen.
\item \emph{Kernaussage (1.7):} man sieht sofort, welche Aufgaben parallel laufen und welche
      zusammenhängen.
\item \emph{Nutzung (1.10):} Plan bleibt in der abgegebenen Form, die Statusspalte dient zum Verfolgen
      des Fortschritts. \textbf{Achtung:} Janosch schreibt, der Plan sei eine \enquote{Arbeitshypothese,
      die wir anpassen} -- mit ihm abstimmen, damit sich eure Abgaben nicht widersprechen.
\end{enumerate}
\end{hinweise}
\begin{umfang}{Einlage 1 (Arbeitsplan): Beschreibung}{0.25}{0.5}
\noindent\emph{\enquote{Dieser Text ist von KI formuliert worden.}}~\cite{Manevski2026KIArbeitsplan}

Bei dieser Einlage handelt es sich um den Arbeitsplan für unser Projekt \enquote{\Projekt}
(Anhang~\ref{anh:arbeitsplan}). \Partner{} und ich haben ihn gemeinsam in Excel als Gantt-Diagramm
aufgebaut, im Wochenraster von KW~42/2026 bis KW~11/2027. Die Baumstruktur, in der jedes parallel laufende
Thema seine eigenen Zweige hat, haben wir uns beim Beispiel \enquote{Mittagessen kochen} aus der
Aufgabenstellung abgeschaut \cite{MoodleArbeitsplan2026}. Fertig war der Plan am 22.\,10.\,2026, am
selben Tag haben wir ihn mit Prof.~Marti besprochen \cite{Marti2026Gespraech}.

Gegliedert ist er in sieben Blöcke, von Theorie und Literatur über Materialien, Versuchsaufbau, Messungen
und Auswertung bis zu Vortrag und Portfolio. Zu jeder Aufgabe steht drin, wer sie übernimmt, wann sie
beginnt und endet und wie weit sie ist; pro Block gibt es eine Prozentanzeige, und alle Abgaben sind
markiert. Die wichtigste Aussage ist für mich, dass man auf einen Blick sieht, welche Aufgaben parallel
laufen und welche zusammenhängen. Als Meilensteine sehe ich den funktionierenden Prototyp, die fertige
Simulation und eine saubere Messreihe.
\end{umfang}

\subsection{Reflexion}
\begin{hinweise}
In der Reflexion sollen Sie sich Gedanken machen über die Wirkung dieser Einlage und der Arbeit mit ihr auf
Ihren Lernfortschritt.
\begin{itemize}\itemsep0pt
\item In welchem Zusammenhang steht der Arbeitsplan zu Ihrem Projektthema?
\item In welchem Zusammenhang steht der Arbeitsplan zu Ihrem Schwerpunktthema?
\item Was haben Sie bei der Erstellung des Arbeitsplanes gelernt?
\item Wie hat Ihnen der Arbeitsplan beim Verständnis Ihres Projektthemas und Ihres Schwerpunktthemas geholfen?
\item Wo gab es Schwierigkeiten bei der Beschäftigung mit der Arbeitsplan?
\item Was sagen Ihre Kommilitoninnen oder Kommilitonen dazu, als Kommentar oder wenn Sie mit Ihnen über den Arbeitsplan diskutieren?
\item Enthält Ihr Arbeitsplan logische oder physikalische Fehler?
\item Gab es bei der Beschäftigung mit dem Arbeitsplan eine überraschende neue Erkenntnis?
\end{itemize}
\textbf{Die Reflexion für eine Einlage sollte mindestens eine Seite und maximal 2 Seiten lang sein. Die
Längenangaben sind ohne Titel, Bilder, Tabellen und Tabellen- oder Bildlegenden zu verstehen.}
Peer-Feedback zitieren (siehe \texttt{literatur.bib}, Vorlage \texttt{VORLAGE-Gespraech}).

\smallskip\emph{Deine Gliederung aus dem Fragebogen (Ziel ca.\ 1{,}2 Seiten, vier Absätze, Ich-Form,
eher sachlich mit persönlichen Stellen):}
\begin{enumerate}\itemsep0pt
\item \emph{Wie der Plan entstanden ist (A1--A4, D1, D3, D4):} anfangs nicht in Worte zu fassen, was
      wir wollten; Brainstorming und gezeichnete \enquote{Prototypen}; mehrere Runden im Kopf und auf
      Papier, erst dann grob in Excel. Zu jedem Schwerpunkt alle Aufgaben gesammelt, aufgeteilt, nach
      Zusammenhang und sinnvollem Zeitabschnitt gruppiert, Dauer diskutiert. Geholfen hat die Erfahrung
      aus zwei Jahren Grundpraktikum: \enquote{schnelle} Aufgaben sind nie schnell erledigt. Am Anfang
      Meinungsverschiedenheiten mit \Partner, gelöst über Argumente und vergessene Themen. Gelernt: wie
      wichtig Iterieren ist -- mit Abstand sieht man, was fehlt, und hinterfragt Erarbeitetes. Noch nie
      einen so detaillierten Plan gemacht.
\item \emph{Projektthema (B1--B4, G7):} vorher nur grob bekannt, wie der Stein funktioniert. Erst beim
      Ableiten der Aufgaben aus dem Versuchsziel wurde klar, was in der Theorie steckt. Erst beim Planen
      als eigene Schritte erkannt: Reibung, reproduzierbarer Anstoß, Sensorkalibrierung. Aufwendiger als
      gedacht: Theorie und Vorbereitung, das 3D-Modell muss für den kleinen Effekt sehr genau sein, und
      schon kleinste Änderungen der Startbedingungen spielen eine große Rolle. Offen: ob Glasplatte und
      Anstoß das überhaupt reproduzierbar erlauben.
\item \emph{Schwerpunkt (C2--C5, H2):} eigene Zeilen: Bewegungsgleichungen (1.4, 29.\,10.--25.\,11.),
      Firmware und Kalibrierung (2.1.3/2.1.4, 05.--26.\,11.), Simulation (1.6, 19.\,11.--17.\,12.),
      Datenpipeline (5.1, 26.\,11.--17.\,12.), Video (3.2). \textbf{Aus dem Plan abgelesen:} vom 19.\ bis
      25.\,11.\ laufen 1.4, 1.6, 2.1.3 und 2.1.4 gleichzeitig -- dein Engpass; die Simulation startet
      am 19.\,11., die Körperparameter von Janosch (1.5) sind erst am 25.\,11.\ fertig -- deine direkte
      Abhängigkeit von ihm (alle wichtigen Größen des Körpers). Motivation: Programmierung, Theorie und
      Simulation zwingen dazu, sich um kleinste Details zu kümmern, die sonst im Verständnis fehlen
      würden. Erst wenig Lust aufs Projekt erwartet, nach dem Ausarbeiten der eigenen Aufgaben jetzt Spaß.
\item \emph{Schwierigkeiten, Rückmeldungen, Ausblick (D2, E1--E3, F2--F4, G, I1, I4):} der Gedanke an
      so einen detaillierten Plan war anfangs einschüchternd, danach eher Zufriedenheit. Neigung, Zeit
      zu unterschätzen -- bekannt und bewusst mit Puffer kompensiert, v.\,a.\ beim 3D-Druck (Teillieferungen)
      und weil wir vom Betreuer abhängen, der den Stein druckt. Prof.\ Marti hat Fragen gestellt und den Plan
      dann so akzeptiert \cite{Marti2026Gespraech}; von Kommilitonen nur, dass er anfangs hässlich war --
      optisch überarbeitet, Struktur gleich gelassen. Auffällige Zeitüberschneidungen geprüft: gewollt als
      Puffer bzw.\ weil die Messungen selbst kurz sind. Kontrolle über ständige Absprache und die
      Statusspalte. Was ich beim nächsten Plan anders mache, kann ich erst beurteilen, wenn ich versucht
      habe, diesen einzuhalten.
\end{enumerate}
\emph{Unabhängig von Janosch bleiben} (Moodle: Reflexionen \enquote{einzeln und unabhängig}): nicht seine
Absatzthemen übernehmen (Messtechnik/Puffer, Messkonzept Messreihe~3, Detailgrad) und nicht seine
Formulierungen (\enquote{vom Ende her denken}, \enquote{Arbeitshypothese}).
\end{hinweise}
\begin{umfang}{Einlage 1 (Arbeitsplan): Reflexion}{1}{2}
\noindent\emph{\enquote{Dieser Text ist von KI formuliert worden.}}~\cite{Manevski2026KIArbeitsplan}

Am Anfang konnten Janosch und ich gar nicht richtig in Worte fassen, wie der Plan aussehen soll. Wir haben
deshalb erst gebrainstormt und ein paar \enquote{Prototypen} auf Papier gezeichnet; der Plan ist also
zuerst im Kopf entstanden, dann in mehreren Runden auf dem Blatt und erst zum Schluss in Excel. Da unsere
Schwerpunkte schon feststanden, haben wir zu jedem Schwerpunkt alle Aufgaben aufgeschrieben, aufgeteilt und
danach sortiert, was zusammenhängt und in welchen Abschnitt des Praktikums es passt. Am längsten haben wir
darüber diskutiert, wie viel Zeit jede Aufgabe bekommt. Geholfen hat mir dabei die Erfahrung aus dem
Grundpraktikum: Aufgaben, bei denen man denkt, dass sie schnell erledigt sind, sind eigentlich nie schnell
erledigt. Anfangs waren
wir uns nicht überall einig, aber mit Argumenten und Themen, die einer von uns vergessen hatte, sind wir bei
einem Plan gelandet, mit dem wir beide zufrieden sind. So einen detaillierten Plan hatte ich vorher noch nie
gemacht. Gelernt habe ich vor allem, wie viel das Iterieren bringt: Schaut man nach einer Weile mit anderem
Blick auf das schon Erarbeitete, merkt man, was fehlt, und fragt sich, ob alles wirklich stimmt und wichtig ist.

Über den Wackelstein wusste ich vorher eigentlich nur, dass es ihn gibt und wie er grob funktioniert. Erst
als wir aus dem Ziel des Versuchs die einzelnen Aufgaben abgeleitet haben, ist mir klar geworden, wie viel
Theorie hinter dem Effekt steckt, die man beim bloßen Anschauen des Steins nicht sieht. Die Reibung auf der
Glasplatte, einen reproduzierbaren Anstoß oder die Kalibrierung des Sensors hatte ich nie als eigene
Arbeitsschritte gesehen, dabei sind genau das die Stellen, an denen der Versuch scheitern kann. Aufwendiger
als gedacht sind vor allem Theorie und Vorbereitung: Für so einen kleinen Effekt muss das 3D-Modell praktisch
perfekt sein, und schon kleinste Änderungen der Startbedingungen spielen eine große Rolle. Ob die Glasplatte
Rollen ohne Gleiten zulässt und ob sich der Anstoß reproduzieren lässt, müssen wir noch herausfinden.

Mein Schwerpunkt sind die Bewegungsgleichungen und die Simulation der Spin-Umkehr. Im Plan gehören dazu die
Bewegungsgleichungen (1.4), Firmware und Kalibrierung des Sensors (2.1.3, 2.1.4), die Simulation (1.6) und
die Datenpipeline in Python (5.1). Beim Einteilen dieser Aufgaben habe ich besser verstanden, was mein
Schwerpunkt konkret bedeutet und womit er zusammenhängt. Vom 19.\ bis 25.\,11.\ laufen zum Beispiel
Bewegungsgleichungen, Simulation, Firmware und Kalibrierung gleichzeitig, das wird meine vollste Woche.
Außerdem brauche ich für die Simulation alle wichtigen Größen des Körpers von Janosch (1.5), die laut Plan
erst am 25.\,11.\ fertig sind, während die Simulation schon am 19.\,11.\ beginnt. An dieser Stelle hängt mein
Teil am stärksten von seinem ab. Gereizt haben mich an diesen Aufgaben das Programmieren, die Theorie und die
Simulation, weil man sich dabei um jedes Detail kümmern muss, an das man ohne dieses Ziel nie denken würde
und das dann im Verständnis fehlt. Ehrlich gesagt hatte ich erwartet, dass mir das Projekt wenig Spaß macht;
seit ich meine Aufgaben ausgearbeitet habe, sehe ich das anders.

Wirklich schwer gefallen ist mir beim Planen nichts, auch wenn mich der Gedanke an so einen detaillierten
Plan am Anfang eher eingeschüchtert hat. Ich weiß von mir, dass ich die Zeit für Aufgaben gern unterschätze.
Das haben wir bewusst mit viel Puffer ausgeglichen, besonders beim 3D-Druck und bei Teillieferungen, wo wir
zusätzlich von unserem Betreuer abhängen, der den Stein druckt. Prof.~Marti hat zu dem Plan einige Fragen
gestellt, die wir beantworten konnten, und ihn danach so akzeptiert \cite{Marti2026Gespraech}. Von
Kommilitonen kam vor allem, dass der Plan anfangs nicht schön aussah; wir haben ihn daraufhin optisch
überarbeitet, an der Struktur aber nichts geändert. Logische Fehler sehe ich keine: Wo sich Zeiträume
überschneiden, ist das gewollt, als Puffer für Nachmessungen oder weil die Messungen selbst nicht lange
dauern. Auf dem Laufenden halten wir uns über ständige Absprache und die Statusspalte. Was ich beim nächsten
Plan anders machen würde, kann ich erst sagen, wenn ich versucht habe, diesen hier einzuhalten.
\end{umfang}

%% Optional (vom Leitfaden ausdrücklich erwünscht): eigene Bewertungskriterien
%\subsection{Eigene Bewertungskriterien}
%Für diese Einlage schlage ich zusätzlich folgende Kriterien vor: ...
